\section{Unit-pivot rays before disc orbit discovery}
\label{sec:unit-ray-native}

Report 69 supplies a particularly small sufficient test for a one-dimensional
normal matching support. The maintained disc-count API now tries that test
after its existing vertex-link removal and quadrilateral-gcd reduction.
Success gives a checked connected primitive surface, whose Euler
characteristic and boundary homology determine the essential-disc count.
Failure uses the complete previous component observer. This is an adaptive
choice of proof method based on an inexpensive structural certificate;
it does not guess connectedness from coefficient size or from cohomological
primitivity. The unit producer and checker are copied unchanged from the
delivery; the surrounding disc protocol and native validation are new.

\subsection{Triangular unit rows certify the integral ray}

Let $x\ne0$ be the validated nonnegative admissible vector, $S$ its positive
support, and $M_S$ the matching matrix restricted to that support. Choose a
seed $s\in S$. A transcript lists $|S|-1$ pairs $(\text{row},\text{pivot})$.
At each step the row has exactly one supported coordinate not yet known,
its pivot, and its coefficient is $+1$ or $-1$. After the final step all
supported coordinates are known. The source already establishes existence;
the transcript establishes uniqueness.

\begin{proposition}[Integral unit-ray witness]
Such a transcript proves $\dim_{\mathbb Q}\ker M_S=1$. There is an integral
kernel generator $u$ with $u_s=1$, and $x=x_su$. In particular
$\gcd_{i\in S}x_i=x_s$.
\end{proposition}
\begin{proof}
Set the seed value to one. Each row solves its new coordinate as an integer
linear combination of earlier known coordinates, since its pivot is a unit.
In discovery order these $|S|-1$ rows have a triangular minor with diagonal
entries $\pm1$, so their rank is $|S|-1$. The nonzero matching source gives
nullity at least one. Thus every kernel vector is determined by its seed,
and the constructed $u$ equals $x/x_s$, satisfies all remaining equations,
and is positive on $S$. Its integral coordinates include one, proving the
gcd assertion. No values need to be propagated in the certificate planner:
the already validated source and the ordered row pattern suffice.
\end{proof}

The producer chooses the minimum-bit positive coordinate, breaking ties by
index. It keeps incident rows, an unknown count and the XOR of unknown
column indices for each row. Processing a known column decrements each
incident count once. A row reaching one unknown exposes that column if its
coefficient is a unit. Each supported column is processed once and each
incidence once. The independent checker instead reconstructs the support
and examines the chosen original rows in order; it imports no planner.
Matching rows have bounded local size, so both use linear sparse-incidence
work, besides geometry preparation and source arithmetic. The indexed
transcript has $O(|S|\log(t+1))$ bits. This is not a linear-bit-time claim
for reading arbitrarily large coordinates.

Failure of this gate does not prove nullity greater than one. For example
the rank-one equation $2x_0-3x_1=0$ with source $(3,2)$ has no unit pivot.
Multiple free variables, an empty support, or a seed that does not expose
all columns also miss this sufficient criterion. No repeated seed search
or Gaussian elimination is introduced on a miss.

\subsection{Connectedness and the exact essential-disc count}

Write the original source as vertex links plus $m v$, where the existing
canonical routine removes maximal complete triangle layers and divides by
the quadrilateral gcd $m$. A nonzero reduced core $v$ is primitive: its
quadrilateral gcd is one and matching forces the peeled triangle content
to divide by the same gcd. A successful unit transcript therefore has
seed coordinate exactly one; both producer and checker enforce it.

\begin{lemma}[Primitive supported ray is connected]
If an admissible primitive normal vector $v$ spans its supported matching
kernel, then its normal surface is connected.
\end{lemma}
\begin{proof}
Each component has a nonnegative integral matching vector supported within
$v$, hence equals $\lambda v$ for some positive rational $\lambda$.
For a unit seed, $\lambda$ is the component's integral seed value. More
generally primitivity and a B\'ezout combination of the entries also force
$\lambda$ to be integral. These positive integers sum to one. There can
be only one component.
\end{proof}

In the compact orientable manifold with torus boundary validated by the
existing geometry, a connected surface with $\chi(v)=1$ is a disc when
its boundary is nonempty; the closed alternative is a projective plane.
The boundary class is computed from two checked torus homology cycles,
by intersecting each with the normal boundary modulo two. A nonzero class
excludes the closed alternative and proves essential disc boundary. Conversely
an essential simple circle on a torus has primitive slope, so its two
mod-two coordinates cannot both vanish. Therefore
\[
 D(mv)=m\,\mathbf1\{\chi(v)=1\text{ and }[\partial v]\ne0
                       \text{ in }H_1(\partial T;\mathbb F_2)\}.
\]
Here $D$ counts essential-disc components, not all components.

The zero case is also exact: any disc component of $mv$ has vector $kv$
for a positive integer $k$, and $1=k\chi(v)$ forces $k=1$ and
$\chi(v)=1$. Its essential boundary supplies the other condition. A disc
is two-sided in this ambient manifold, so a positive result consists of
$m$ parallel copies. One-sided primitive rays need not have $m$ total
components: even scaling gives their connected two-sided frontier. That
frontier has Euler characteristic $2\chi(v)$ and cannot create a disc.
Removed vertex links are spheres or boundary-parallel discs and contribute
zero to $D$. This proves the displayed count for the whole original source.

The argument uses rank and integral coordinates. The disconnected
primitive-cohomology example of Section~\ref{sec:completion-theorems} does
not satisfy this reasoning merely because its class is primitive.

\subsection{The whole Fibonacci family has short witnesses}

The preserved report proves more than the finite tests. In its actual
layered fixture, the three supported coordinates per tetrahedron are
$a_i=b_i=F_{t-i+1}$ and $c_i=F_{t-i}$, with $F_1=F_2=1$.
The original face equations for $i<t-1$ include
\[
 b_i=b_{i+1}+c_{i+1},\quad c_i=a_{i+1},\quad
 a_i=a_{i+1}+c_{i+1},\quad c_i=b_{i+1},
\]
and the terminal triple is equal. For $t=1$ the chosen seed is $a_0$;
otherwise it is $c_{t-2}$. The equalities expose the terminal triple,
then the recurrence exposes each previous triple. Every pivot is a unit,
so there are exactly $3t-1$ discovery steps for every $t\ge1$ and positive
scale $m$. The next larger value is at least twice the seed, keeping the
minimum-bit/index choice fixed under scaling.

For completeness the source cell counts also prove the topology, rather
than assuming a descriptive fixture name. Its global edge weights are
$F_2,\ldots,F_{t+3}$, giving $V=F_{t+5}-2$; its pieces number
$P=F_{t+5}-5$. The total local side count is
$6F_{t+3}+4F_{t+2}-16$, with $2F_{t+3}$ unpaired boundary arcs, so
$E=2F_{t+5}-8$ and $\chi=V-E+P=1$. Boundary weights on the one-vertex
torus are $F_{t+1},F_{t+2},F_{t+3}$. Adjacent Fibonacci numbers are coprime;
their mod-two cocycle is nonzero. A one-vertex triangulation has no nonzero
vertex coboundaries, so the boundary class is nonzero. The ray theorem
gives exactly $m$ essential discs.

The primitive piece count is exponential in $t$, but the indexed rank
witness has $O(t\log(t+1))$ bits. The full binary source has
$\Theta(t^2+t\log(m+1))$ bits. This is a family theorem, not a linear-time
claim in the tetrahedron count for the entire call. Unlike the delivery's
full ray-block API, the native adaptation deliberately retains its existing
gcd and division prefix. Only the subsequent unit-witness stage avoids
value propagation and elimination. The whole call is still polynomial in
its encoded supplied input; its proof and orbit work are substantially
smaller on these dense rays.

\subsection{Protocol, fallback and resource semantics}

\code{normal\_compressing\_disk\_count} defaults to
\code{unit\_ray=True}, with strict Boolean option validation. A successful
gate returns \code{normal-disc-count-v2}; its seven outer fields bind the
input fingerprint, divisor, vertex links, core coordinates, core certificate
and final count. The inner \code{normal-unit-ray-disc-v1} has exactly six
fields: its schema, support transcript, torus homology basis, Euler value,
two mod-two boundary coordinates and primitive disc count. The transcript
itself has strict support, seed, step and nullity fields.

Replay independently rebuilds the original geometry, source binding,
link/gcd reduction, matching support, triangular rank evidence, torus basis,
Euler characteristic and boundary parities. It invokes neither the planner,
the canonical reduction producer nor orbit discovery. Shared finite geometry
and arithmetic helpers remain part of the trusted mathematical boundary;
the producer and replay control paths are distinct. An incorrect primitive
seed is rejected, rather than used as an asserted multiplicity.

An unsuccessful gate retains the complete component query and its
version-one proof. Setting \code{unit\_ray=False} reproduces the old
schedule and result exactly.
Pure vertex links retain the old zero-core proof. Existing version-one
certificates remain accepted. The historical coherent-obstruction driver
explicitly disables the gate to preserve its published proof bytes.

The shortcut reports zero orbit cycles, zero weight runs and zero replay
events, with its own ray-step count. A \code{max\_cycles=0} allowance bounds
orbit discovery, so it permits this algebraic path. The supplied-sector
test now correctly finds a certified meridian with zero orbit allowance;
a zero basis-enumeration allowance still prevents candidate discovery.
Callback cancellation covers preparation, planning and replay and propagates
to callers. A missed gate with an exhausted orbit allowance remains
inconclusive; it does not certify a negative knot verdict.

The source pipeline still matters. The supplied-sector producer benefits
through its disc checker, but the main cocycle candidate schedule does not
use this API and has not acquired new source coverage. This integration is
an exact observer improvement for a supplied vector or a supplied sector,
not a new complete diagram-recognition implementation.

\subsection{Validation and complete-call timings}

All 1,317 maintained test methods pass in 236.139 seconds. The focused
34-method run passes in 2.896 seconds. New tests cover layered sizes
1, 2, 8, 32, 128 and 256; exact $3t-1$ transcripts; relabelings; links and
binary scalings through $2^{20,000}$; one-sided negatives; nonunit rank-one
misses; strict option types; callback cancellation; and field/row/seed
mutations with producer-disabled replay.

The corpus audit compares 1,275 vectors under both periodic rules, or
2,550 configurations, against the frozen \code{66098968e} count producer
and checker and the corpus component oracle. Every count agrees. The gate
succeeds on 1,878 configurations and uses zero orbit cycles; the 672 misses
reproduce the complete baseline result. Disabling the gate reproduces all
2,550 baseline results exactly. The audit retains 44 representative complete
proofs and 553 source hashes. Verification of old and new proof
languages is checked separately. The corpus contains duplicates by scale
and relabeling, so 1,878 is not a count of distinct knot types.

The isolated benchmark uses five rounds and four interleaved arms,
including old/old and new/new controls: 220 measured calls and 44 warmups
all complete. Each timed call includes production, independent replay and
serialization/hash of the complete certificate. The count reference is
the frozen producer \emph{and} checker; for supplied-sector calls its disc
producer is replaced by that baseline and the maintained outer replay
accepts its legacy proof. All arms agree on counts or discovered-coordinate
hashes, and proof bytes are deterministic within each arm. Timing ratios
below are medians of paired old/new ratios, not ratios of separate medians.

\noindent\textbf{Supplied-vector disc counts.}
\begin{center}
\small
\begin{tabular}{lrrrrr}
\toprule
Input & old ms & new ms & old/new & old A/A & new A/A\\
\midrule
Layered 16 & 12.891 & 3.556 & 3.639 & 0.964 & 1.014\\
Layered 64 & 229.634 & 26.748 & 8.685 & 1.010 & 0.997\\
Layered 256 & 1811.970 & 57.238 & 31.386 & 1.003 & 1.012\\
Layered 64, binary scale & 130.841 & 19.010 & 6.924 & 1.036 & 0.997\\
Empty & 0.226 & 0.226 & 0.992 & 0.988 & 1.000\\
One-sided & 0.777 & 0.390 & 1.995 & 0.998 & 0.989\\
Vertex link & 0.235 & 0.231 & 1.019 & 1.007 & 1.003\\
Sphere/disc/Mobius mix & 1.898 & 1.010 & 1.901 & 0.994 & 0.994\\
Sphere/negative Euler mix & 4.298 & 1.616 & 2.654 & 1.001 & 0.995\\
\bottomrule
\end{tabular}
\end{center}

\noindent\textbf{Complete supplied-sector discovery.}
\begin{center}
\small
\begin{tabular}{lrrrrr}
\toprule
Input & old ms & new ms & old/new & old A/A & new A/A\\
\midrule
Layered 16 & 55.591 & 42.830 & 1.303 & 1.001 & 1.002\\
Layered 64 & 1945.525 & 1802.229 & 1.068 & 1.006 & 1.013\\
\bottomrule
\end{tabular}
\end{center}


The layered-256 query drops from about 1.812 seconds to 0.057 seconds,
a paired 31.39-fold gain. Its complete certificate shrinks from 7,098,789
to 36,240 bytes. Layered-64 improves 8.68-fold; its $2^{4096}+1$ multiple
improves 6.92-fold after paying the existing reduction prefix. Empty and
pure-link controls are near parity. The mixture cases benefit because
link removal leaves a certified ray; this is not promotion of a generic
multi-block decomposition.

Complete supplied-sector discovery improves only 1.30-fold at size 16 and
1.07-fold at size 64 because enumeration remains dominant. The A/A medians
range from 0.964 to 1.036; the raw records preserve all rounds and bytes.
No full knot-recognition timing is inferred from these measurements.
Reproduction is \code{python -B -m normal\_orbit\_research.unit\_rays audit}
and \code{benchmark --rounds 5}, each with a required \code{--output}
path, from \code{fast/}. The retained records
are \code{data/unit-ray-audit.json}, \code{data/unit-ray-benchmark.json}
and both test logs. The remaining global problem is controlled complete
witness discovery, not the already polynomial inspection of one supplied
normal vector. No general quasipolynomial bound follows from this shortcut.
